\subsection[Products with arbitrary effective C1 residue fields]{Products with arbitrary effective $C_1$ residue fields}
\label{C:C1-coherent-products-subsection}

\begin{theorem}[Proved: bounded product coordinates with full square classes]
\label{C:C1-coherent-products-WG}
For each $i$ let $R_i$ be one of the solved local
rings of Theorems~\ref{C:effective-henselian-DVR-WG},
\ref{C:finite-field-henselization-WG},
\ref{C:arithmetic-henselization-WG},
\ref{C:arithmetic-strict-henselization-WG},
\ref{C:effective-henselian-flat-unions},
\ref{C:C1-henselization-WG}, or
\ref{C:effective-coherent-henselization-WG}.
Thus every factor has literal canonical indices,
actual generators and ordinary terminating local
classification and operation algorithms.
Put $R=\prod_iR_i$.

At a symmetric indecomposable index $j$ of
one factor write its coordinate as
$(a_j,\gamma_j)\in
\mathbf Z\oplus K_j^\times/(K_j^\times)^2$,
where $\gamma_j$ has its proved canonical
representative.  For a finite residue field
this is the earlier determinant bit, and
for algebraically closed residue
$\gamma_j=1$.
All other types have an integer coordinate
$a_j$.  Each factor's coordinate list has
finite support.
Let $n_j,d_j$ be its canonical module's
rank and minimal coefficient degree, and put
\[
 w_j=\begin{cases}
 n_j,&j\text{ symmetric or unitary},\\
 2n_j,&j\text{ alternating or paired},
 \end{cases}
 \qquad
 \delta_j=\begin{cases}
 1,&j\text{ symmetric},\ a_j=0,\ \gamma_j\ne1,\\
 0,&\text{otherwise}.
 \end{cases}                                             \tag{C.3.CP1}
\]
Define, in each factor,
\[
 \begin{split}
 r_+(\xi)&=\sum_jw_j\max(a_j,0)
                +\sum_{j\text{ symmetric}}n_j\delta_j,\\
 r_-(\xi)&=\sum_jw_j\max(-a_j,0)
                +\sum_{j\text{ symmetric}}n_j\delta_j,\\
 c(\xi)&=\max\bigl(\{d_j:a_j\ne0\text{ or }
                   (j\text{ symmetric},\gamma_j\ne1)\}\cup\{0\}\bigr).
 \end{split}                                             \tag{C.3.CP2}
\]
These are computable simultaneous minima of the
positive rank, negative rank and common degree
of a difference of actual forms representing
$\xi$.

The canonical map
\[
 WG_{\rm sp}(BO_{m,R})\longrightarrow
                    \prod_iWG_{\rm sp}(BO_{m,R_i})       \tag{C.3.CP3}
\]
is injective.  Its image is exactly the families
for which $r_+$, $r_-$ and $c$ are uniformly
bounded on the index set of factors.
Products and every exterior operation are
coordinatewise, and each requested coordinate
is obtained by its ordinary local algorithm.
This is an exact description also when the
square-class groups are infinite.  No finite
decision procedure for equality or boundedness
of an arbitrary infinitely supplied family
is asserted.
\end{theorem}
\begin{proof}
The complete local form classifications give,
in each symmetric type, the monoid of pairs
$(r,\gamma)$ with $r\geq0$ and
$\gamma\in K_j^\times/(K_j^\times)^2$,
subject to $\gamma=1$ for $r=0$.
Orthogonal sum adds ranks and multiplies
square classes.  A positive rank permits
every square class: use a diagonal residue
form with that scalar once and all other
entries one, and lift its actual generator.
This statement applies to the finite
and algebraically closed residue cases
as its specializations.
Each other type has monoid $\mathbf N$
with generator rank $w_j$.
All these monoids are cancellative.

For a symmetric coordinate $(a,\gamma)$,
let $u,v$ be its positive and negative
multiplicities.  Then $u-v=a$ and
their determinant square classes have
quotient $\gamma$.
If $a>0$ the simultaneous minimum is
$(u,v)=(a,0)$ and the positive form
can carry $\gamma$.
If $a<0$ it is $(0,-a)$; the negative
form can carry $\gamma^{-1}=\gamma$,
since every square class has order
dividing two.
If $a=0$ and $\gamma\ne1$, both
multiplicities must be positive and
their minimum is $(1,1)$, represented
by the scalar form minus its reference
form.  The zero coordinate has minimum
$(0,0)$.
The other types have minima
$\max(a_j,0)$ and $\max(-a_j,0)$.
The local classifications make
different blocks independent, so
summing their ranks proves the
first two formulas and constructs
a pair attaining both bounds.

If an actual representation has
coefficients in $C_D$, every direct
summand has degree at most $D$:
over the local factor choose free
bases of the summand and its
complement, then conjugate by
the resulting constant invertible
matrix.  All resulting coefficients
remain in the same $C_D$.
Every supported coordinate forces
its indecomposable or its dual
pair to occur in at least one
of the two forms.
The antipode preserves minimum
degree, so paired duals have
the same $d_j$.
This proves the lower bound
$c(\xi)$ in (C.3.CP2).
The explicit forms attaining
the rank minima use only those
supported indecomposables;
changing their Gram forms
does not change their coactions.
They attain the same degree
bound simultaneously.
The zero class uses two
zero modules of degree zero.
All ranks, degrees and
square-class identity tests
are finite computable data.

A difference of two actual
forms over $R$ has bounded
ambient idempotent sizes
and finite coefficient degree,
and therefore has the three
asserted bounds in every
factor.
If its image in (C.3.CP3)
is zero, local cancellation
gives actual local isometries.
For presentations $P=eR^a$
and $Q=fR^b$, these isometries
and their inverses have
fixed matrix sizes $b\times a$
and $a\times b$, with the
corresponding corner conditions.
Assemble their entries in
the product ring.
All Gram and inverse equations
hold coordinatewise.
The equivariance equations
hold as actual coefficient
equations because every
fixed $C_D$ has a finite
free integral basis, so
its map to the product
of factor coefficient
modules is injective.
The assembled matrices
are therefore an actual
equivariant isometry.
This proves injectivity.

Conversely, choose common
bounds $B,D$ for an input
coordinate family.
Construct its actual local
positive and negative forms
with the simultaneous minima.
Pad each free presentation
to size $B$ with the
idempotent $\operatorname{diag}(I_r,0)$,
and pad the coaction, Gram
matrix and inverse Gram
matrix with zeros on the
complement.
The padded coaction has
that idempotent as its
counit.
Assemble the finitely many
matrix positions and the
fixed finitely many coefficient
positions in $C_D$ across
all factors.
The idempotent, coaction,
form and inverse identities
are finite identities and
hold coordinatewise, hence
over $R$.
They define actual global
positive and negative forms
with exactly the given
coordinate difference.
This proves the image assertion.

Tensor and exterior constructions
commute with base change.
Products of differences with
bounds $(B,D)$ and $(B',D')$
have actual rank bound $2BB'$
and degree bound $D+D'$.
For an exterior coefficient
of order $r$, expand
$\lambda_t(P)\lambda_t(Q)^{-1}$
as a finite integral sum
of tensor products of
exterior forms.
Every term has degree at
most $rD$.
For an explicit rank bound
put $A_0(B)=1$ and
\[
 A_r(B)=\binom Br+
       \sum_{j=1}^r\binom Bj A_{r-j}(B),
 \qquad \binom Bj=0\text{ for }j>B.                     \tag{C.3.CP4}
\]
The formal series
$\sum_r A_r(B)t^r=(1+t)^B/(2-(1+t)^B)$
bounds the sum of the ranks
of all terms before signs
are grouped: its denominator
expands the inverse exterior
series using absolute
coefficient bounds.
Thus each of the positive
and negative output ranks
is at most $A_r(B)$.
This proves that all
operations preserve the
image bounds.
The local ordinary algorithms
compute every requested
output coordinate, which
is precisely the algorithmic
access asserted.
\end{proof}
